\documentclass[10pt,a4paper]{amsart}
\usepackage[margin=19mm]{geometry}
\usepackage{amsmath,amssymb,mathtools}
\usepackage[scr=rsfs,cal=euler]{mathalfa}
\usepackage{xfrac}
\usepackage[unicode,hidelinks,bookmarks=false]{hyperref}
\newcommand{\ie}{i.e.\ }
\newcommand{\cf}{cf.\ }
\newcommand{\resp}{respectively\ }
\newcommand{\Subsection}{\subsection}
\providecommand{\nobreakdash}{\nobreak\hbox{-}}
\setlength{\emergencystretch}{2em}
\multlinegap0pt
\usepackage{mathtools}
\usepackage{enumerate}
\usepackage{amssymb}
\newtheorem{theorem}{Theorem}
\newcommand{\Sig}{\Sigma}
\newcommand{\bbm}{\begin{bmatrix}}
\newcommand{\be}{\begin{equation}}
\newcommand{\bnum}{\begin{enumerate}}
\newcommand{\ddd}{,\ldots,}
\newcommand{\diag}{\mbox{diag}}
\newcommand{\dt}{\delta}
\newcommand{\ebm}{\end{bmatrix}}
\newcommand{\ee}{\end{equation}}
\newcommand{\enum}{\end{enumerate}}
\def\rank{\mbox{rank}}
\newcommand{\re}{\mathbb{R}}
\newcommand{\reff}[1]{(\ref{#1})}
\newcommand{\sig}{\sigma}
\newcommand{\tsig}{\tilde{\sigma}}
\newcommand{\tu}{\tilde{u}}
\newcommand{\tv}{\tilde{v}}
\title{Rank One Completion for Higher Order Tensors}
\author{Linghao Zhang, Ioana Dumitriu, Jiawang Nie}
\date{Source: arXiv:2604.23104v1 (CC BY 4.0)}
\begin{document}
\maketitle

\noindent\emph{Source and licence.} Rank One Completion for Higher Order Tensors. Version arXiv:2604.23104v1 (CC BY 4.0), distributed by arXiv under the Creative Commons Attribution 4.0 International licence, http://creativecommons.org/licenses/by/4.0/, as recorded in the arXiv licence metadata for this exact version and shown on the abstract page https://arxiv.org/abs/2604.23104v1. Full licence notice: \texttt{licenses/arxiv-2604.23104-CC-BY-4.0.txt}.\par\smallskip

\noindent\textit{Editorial scope.} Counted unit: the article's theorem thm:mx\_null\_perturb (source0.tex lines 2590--2606). The article's complete original proof of it is delivered with it (source0.tex lines 2608--2772, one \texttt{proof} environment of 165 lines). No mathematics is written, repaired or re-derived: the statement and the proof are the article's own lines, in the article's own notation, and no retained line is substituted or re-flowed.  No further source-local context is needed: every cross-reference of every retained line resolves inside the two delivered spans.  Nothing was omitted from any retained span, so the package carries no omission record.  No retained line contains a citation, so the wrapper carries no bibliography and no bibliography entry.  No retained line contains a figure environment, an image-inclusion command or drawing code, so no image is delivered and the package declares no asset dependency.  Editorial formatting only, in addition: an AMS \texttt{amsart} preamble that restates the article's own declarations and macros used by the retained lines, namely the environments \texttt{theorem} on separate counters, together with the standard packages those lines need.  The retained environments print in this document as Theorem 1 in order of appearance, whereas the article prints the same items with the numbering of its own full document.  The complete native source of the article as distributed in the arXiv e-print for this exact version is bundled unchanged as \texttt{sources/199140/source0.tex}.
\begin{theorem}\label{thm:mx_null_perturb}
Let $B \in \re^{\ell \times n}$ be a matrix with $\rank (B) = n-1$.
Let $v$ be a solution to $Bx = 0$ with $\|v\|=1$.
Let $\widetilde{B} = B + \dt B$ be a perturbation of $B$.

\bnum
\item[(i)] If $\ell \ge n$, let $\tv$ be the right singular vector of $\widetilde{B}$ corresponding to the minimum singular value $\tsig_n$ such that $\|\tv\|=1$ and $v^T\tv \ge 0$.

\item[(ii)] If $\ell<n$, let $\tv$ be a solution to $\widetilde{B}x = 0$ with $\|\tv\|=1$ and $v^T\tv \ge 0$.
\enum
Let $\dt v \coloneqq \tv - v$.
Then there exists a constant $C_0>0$, depending on $B$, such that
\be\label{dv<=C0dB}
\|\dt v\| \,\le\, C_0 \|\dt B\|,
\ee
when $\|\dt B\|$ is sufficiently small.
\end{theorem}

\begin{proof}
Let $B = U\Sig V^T$ be the SVD with 
\[
\Sig = 
\bbm 
\Sig_1 & 0 \\
0 & 0
\ebm
\in \re^{\ell \times n},
\quad \Sig_1 = \diag(\sig_1 \ddd \sig_{n-1}),
\]
where $\sig_1 \ge \cdots \ge \sig_{n-1} > \sig_n = 0$ are singular values of $B$, and
\[
U = \bbm u_1 & u_2 & \cdots & u_{\ell}  \ebm \in \re^{\ell \times \ell},
\quad
V = \bbm v_1 & v_2 & \cdots & v_n  \ebm \in \re^{n \times n}
\]
are orthogonal matrices.
Since $\rank(B) = n-1$, the nullspace of $B$ is one-dimensional,
so the vector $v_n$ is unique up to sign.
We can generally assume $v_n = v$.


\noindent (i) For the case that $\ell \ge n$, 
up to orthogonal transformation, we can assume $B=\Sig$ and $v=e_n$.
Let $\tu$ be the unit length left singular vector of $\widetilde{B}$, corresponding to its minimum singular value $\tsig_n$,
i.e., $\widetilde{B} \tv = \tsig_n \tu$.
Note that
\be\label{tBtvn}
\widetilde{B} \tv = (B +\dt B) (v + \dt v) = Bv + B \cdot \dt v + \dt B \cdot v + \dt B \cdot \dt v.
\ee
Let $y \coloneqq B\cdot \dt v$.
Since $Bv = 0$ and $\widetilde{B} \tv = \tsig_n \tu$, the above implies
\[
y = \tsig_n \tu - \dt B \cdot v - \dt B \cdot \dt v.
\]
Write
$
y = 
[ y_1 \,\, y_2 \,\, \cdots \,\, y_{\ell} ]^T.
$
Since $\|\tu\| = \|v\| = \|\tv\| = 1$ and $\|\dt v\| = \|\tv - v\| \le 2$, 
\[
\|y\| \le \tsig_n \|\tu\| + \|\dt B\| \|v\| + \|\dt B\| \|\tv - v\| \le \tsig_n + 3\|\dt B\|.  
\]
Note $\tsig_n$ (resp., $\sig_n$) is the minimum singular value of $\widetilde{B}$ (resp., $B$),
so
\begin{align*}
\tsig_n &= \min_{\|x\|=1} \|\widetilde{B}x\| \le \min_{\|x\|=1}( \|Bx\| + \|\dt B \cdot x\| ) \\
&\le \min_{\|x\|=1} \|Bx\| + \|\dt B\| = \sig_n + \|\dt B\|.
\end{align*}
Since $\rank(B) = n-1$, we have $\sig_n = 0$.
This implies
$
\tsig_n \le \|\dt B\|
$,
so $\|y\| \le 4\|\dt B\|$.
Denote
$[ z_1 \,\, z_2 \,\, \cdots \,\, z_n ] = \dt v^T$.
Since
$
y = B\cdot \dt v = \Sig \cdot \dt v,
$
we get $z_i = y_i / \sig_i$ for $i \in [n-1]$.
Let 
$
\hat{z} = [z_1  \,\, \cdots \,\, z_{n-1} ]^T,
$
then
\[
\|\hat{z}\|^2 = \sum_{i=1}^{n-1} \frac{y_i^2}{\sig_i^2} 
\le \frac{1}{\sig_{n-1}^2} \sum_{i=1}^{n-1} y_i^2
\le \frac{\|y\|^2}{\sig_{n-1}^2}.
\]
Since
$
\tv = v + \dt v = e_n + \dt v = [ z_1  \,\, \cdots \,\, z_{n-1} \,\, 1 + z_n ]^T$,
we have
\[
1 = \|\tv\|^2 = \|\hat{z}\|^2 + (1+z_n)^2,
\quad 1+z_n = \tv^Te_n = \tv^Tv \ge 0.
\]
Hence, $z_n \le 0$ and
\be\label{z_n<=||zhat||^2}
z_n = \sqrt{ 1 - \|\hat{z}\|^2 } - 1 = -\|\hat{z}\|^2 \Big/ \big(\sqrt{ 1 - \|\hat{z}\|^2 } + 1\big).
\ee
This implies $|z_n| \le \|\hat{z}\|^2$, so
\[
\|\dt v\|^2 = \|\hat{z}\|^2 + z_n^2 \le \|\hat{z}\|^2 + \|\hat{z}\|^4 \le
\frac{\|y\|^2}{\sig_{n-1}^2} \Big( 1 + \frac{\|y\|^2}{\sig_{n-1}^2} \Big).
\]
Recall that $\|y\| \le 4\|\dt B\|$.
If $\|\dt B\| \le \frac{\sig_{n-1}}{4}$, then
\[
\|y\| \,\le\, \sig_{n-1} 
\quad \implies \quad
1 + \|y\|^2 \big/ \sig_{n-1}^2 \,\le\, 2,
\]
so
$\|\dt v\| \le \frac{4\sqrt{2}}{\sig_{n-1}} \|\dt B\|$.
Hence, \reff{dv<=C0dB} holds for 
$C_0 =  \frac{4\sqrt{2}}{\sig_{n-1}}$ 
if $\|\dt B\| \le \frac{\sig_{n-1}}{4}$.





\noindent (ii) For the case that $\ell < n$, we have $\ell = n-1$ since $\rank(B) = n-1$.
In the SVD $B = U\Sig V^T$, we have 
$
\Sig = 
[ \Sig_1 \,\, 0 ] \in \re^{(n-1) \times n}.
$
Since $\widetilde{B}\tv = 0$ with $\|\tv\|=1$, \reff{tBtvn} gives
\[
U\Sig V^T \dt v = B \cdot \dt v = - \dt B \cdot v - \dt B \cdot \dt v.
\]
Multiplying $U^T$ on both sides, we have
\[
\Sig V^T  \dt v = - U^T  \dt B \cdot v - U^T \dt B \cdot \dt v.
\]
Let 
$
y = [ y_1 \,\, y_2 \,\, \cdots \,\, y_{n-1} ]^T = \Sig V^T \dt v.
$
Since $U$ is orthogonal and $\|v\| = \|\tv\| = 1$,
\[
\|y\| \le \|U^T\| \|\dt B\| \|v\| + \|U^T\| \|\dt B\| \|\tv - v\| \le 3\|\dt B\|.
\]
Let $z = [ z_1 \,\, z_2 \,\, \cdots \,\, z_n ]^T = V^T  \dt v$
and
$
\hat{z} = [ z_1 \,\, \cdots \,\, z_{n-1} ]^T,
$
then $y = \Sig z$ implies
\[
\|\hat{z}\|^2 \,\le\, \|y\|^2 \big/ \sig_{n-1}^2.
\]
Since
$
V^T\tv = V^T(v+\dt v) = e_n + z = [ z_1 \,\, \cdots \,\, z_{n-1} \,\, 1 + z_n ]^T,
$
we have
\[
1 = \|\tv\|^2 = \|V^T\tv\|^2 = \|\hat{z}\|^2 + (1+z_n)^2,
\]
\[
1+z_n = (V^T\tv)^Te_n = \tv^TVe_n = \tv^Tv \ge 0.
\]
Similar to \reff{z_n<=||zhat||^2}, one can show $|z_n| \le \|\hat{z}\|^2$ and
\[
\|\dt v\|^2 = \|V^T \dt v\|^2 =
\|\hat{z}\|^2 + z_n^2 
\le \frac{\|y\|^2}{\sig_{n-1}^2} \Big( 1 + \frac{\|y\|^2}{\sig_{n-1}^2} \Big).
\]
Recall that $\|y\| \le 3\|\dt B\|$.
If $\|\dt B\| \le \frac{\sig_{n-1}}{3}$, then
$
\|\dt v\| \le \frac{3\sqrt{2}}{\sig_{n-1}} \|\dt B\|.
$
Hence, \reff{dv<=C0dB} holds for 
$C_0 =  \frac{3\sqrt{2}}{\sig_{n-1}}$ 
if $\|\dt B\| \le \frac{\sig_{n-1}}{3}$.
\end{proof}

\end{document}
